\documentclass[11pt]{article}
% Online appendix of the paper (separate PDF on GitHub). Build: python paper/build.py (after main).
% Shared preamble of main.tex (paper) and supplement.tex (supplementary material).
\usepackage[T1]{fontenc}
\usepackage[utf8]{inputenc}
\usepackage{lmodern}
\usepackage[margin=1in]{geometry}
\usepackage{microtype}
\usepackage{amsmath,amssymb}
\usepackage{graphicx}
\usepackage{booktabs}
\usepackage{tabularx}
\usepackage{multirow}
\usepackage{longtable}
\usepackage{array}
\usepackage{xcolor}
\usepackage[most]{tcolorbox}
\usepackage{enumitem}
\usepackage[font=small,labelfont=bf]{caption}
\usepackage{subcaption}
\usepackage[numbers,sort&compress]{natbib}
\usepackage{CJKutf8}
\usepackage{xurl}
\usepackage[section]{placeins}
\usepackage[bookmarks=true,colorlinks=true,linkcolor=blue!50!black,citecolor=blue!50!black,urlcolor=blue!50!black]{hyperref}
\hypersetup{pdftitle={Can Typed Decision Models Serve as Automated Safeguards in PoL Governance? A Clause-Level Analysis of the Proof of Love2 (PoL2) Specification Against Early Evidence},pdfauthor={[Author Name]}}
\usepackage[capitalise,noabbrev]{cleveref}
% Cross-references between the paper and its supplement (labels imported by build.py)
\crefname{supplement}{Supplement}{Supplements}
\Crefname{supplement}{Supplement}{Supplements}

\definecolor{f4pblue}{HTML}{0F4D92}
\definecolor{f4pred}{HTML}{B64342}
\definecolor{f4pgreen}{HTML}{8BCF8B}

\newcommand{\zh}[1]{\begin{CJK}{UTF8}{gbsn}#1\end{CJK}}
\newcommand{\axis}[1]{\textbf{#1}}
\newcommand{\jev}{\textsc{Jev}}
\newcommand{\pol}{PoL2}
\newcommand{\clause}[1]{\S\,#1}
\newcommand{\Clause}[1]{\S\,#1}
\newcommand{\art}[1]{Art.~#1}
\newcommand{\Art}[1]{Article~#1}
% Source marks: Zenodo records (not moderated by a preprint server); grey literature (blogs, vendor pages, repositories)
\newcommand{\zmark}{\textsuperscript{$\ddagger$}}
\newcommand{\gmark}{\textsuperscript{$\dagger$}}
% Corpus counts (generated by make_counts.py; update together with the data files)
% Generated by scripts/merge_coding.py; do not edit by hand.
\newcommand{\nArxiv}{79}
\newcommand{\nZenodo}{27}
\newcommand{\nOther}{28}
\newcommand{\nPapers}{107}
\newcommand{\nGrey}{7}
\newcommand{\nCoded}{85}
\newcommand{\nRows}{161}
\newcommand{\nS}{16}
\newcommand{\nQ}{97}
\newcommand{\nN}{48}
\newcommand{\nPolClauses}{28}
\newcommand{\nKappaN}{30}
\newcommand{\nKappaAgree}{26 of 30}
\newcommand{\nKappa}{0.72}
\newcommand{\nKappaLo}{0.42}
\newcommand{\nKappaHi}{0.94}
\newcommand{\nKappaMarg}{S 3, Q 23, N 4}
\newcommand{\nStrongS}{7}
\newcommand{\nStrongQ}{45}
\newcommand{\nStrongN}{27}
\newcommand{\nStrongRows}{79}
\newcommand{\nNoZenS}{13}
\newcommand{\nNoZenQ}{75}
\newcommand{\nNoZenN}{29}
\newcommand{\nNoZenRows}{117}
\newcommand{\nStronger}{37}
\newcommand{\nModerate}{17}
\newcommand{\nWeaker}{38}
\newcommand{\nAppraised}{92}

% Update-search counts (8 October 2026; generated by scripts/merge_update.py)
% Generated by scripts/merge_update.py (2026-10-08 update search); do not edit by hand.
\newcommand{\nUDocs}{57}
\newcommand{\nUCoded}{48}
\newcommand{\nUArxiv}{27}
\newcommand{\nUZenodo}{21}
\newcommand{\nURows}{118}
\newcommand{\nUS}{12}
\newcommand{\nUQ}{92}
\newcommand{\nUN}{14}
\newcommand{\nUStrongRows}{59}
\newcommand{\nUStrongS}{5}
\newcommand{\nUStrongQ}{45}
\newcommand{\nUStrongN}{9}
\newcommand{\nUKappaN}{118}
\newcommand{\nUKappaAgree}{97 of 118}
\newcommand{\nUKappa}{0.55}
\newcommand{\nUKappaLo}{0.36}
\newcommand{\nUKappaHi}{0.70}
\newcommand{\nUSetPostWindow}{30}
\newcommand{\nUSetLateIndexed}{5}
\newcommand{\nUSetReScreened}{13}

\newcolumntype{L}{>{\raggedright\arraybackslash}X}

\setlist{itemsep=2pt,topsep=4pt}
\newtcolorbox{answerbox}{breakable,enhanced,colback=f4pblue!5,colframe=f4pblue!60,boxrule=0.6pt,
  arc=1.5pt,left=6pt,right=6pt,top=4pt,bottom=4pt,before skip=6pt,after skip=8pt,
  title={\small\bfseries Answer in brief},fonttitle=\small\bfseries,coltitle=white,colbacktitle=f4pblue!80}


% Labels of the paper (written by build.py), so that \cref prints the paper's numbers
\InputIfFileExists{main_labels}{}{}
\renewcommand{\thesection}{\Alph{section}}
\counterwithin*{table}{section}
\renewcommand{\thetable}{\thesection\arabic{table}}
\hypersetup{pdftitle={Online Appendix: Can Typed Decision Models Serve as Automated Safeguards in PoL Governance?}}
\def\UrlBreaks{\do\/\do\-\do\_\do\.}
\title{\sffamily\bfseries Online Appendix\\[4pt]\large Can Typed Decision Models Serve as Automated Safeguards in PoL Governance?\\ A Clause-Level Analysis of the Proof of Love2 Specification}
\author{Shentao Yan}
\date{}

\begin{document}
\maketitle
\noindent Section, table and figure numbers without a letter refer to the paper. The coding data (one row per study and requirement, with code, comparator, design rating, rationale and locator), the catalogue of studies, the excluded records and the scripts are in the \nolinkurl{data/} and \nolinkurl{scripts/} directories of the same repository.

\section{Clause concordance}
\label{oa:concordance}

\Cref{tab:concordance} quotes every \pol{} passage that this paper cites or paraphrases, in the original Chinese~\cite{pol2text} and in the official English translation~\cite{pol2en}, both taken verbatim from commit \texttt{5791ae3}; ``\ldots'' and ``……'' mark omitted text, including where separate list items are joined, and quotation marks inside quotations have been normalised to typographic quotes.
Where a clause has numbered items, the item is given in parentheses.
The right-hand column gives the number used in the snapshot \texttt{2e094d3} cited by the earlier dataset release~\cite{survey01}, so that the two versions can be read side by side; that snapshot numbered \clause{4.3.3.3} as ``3.3.3'' in its heading, which we reproduce.
Both texts are released under CC0~1.0.

\begingroup
\small
\renewcommand{\arraystretch}{1.2}
\begin{longtable}{@{}p{0.085\linewidth}p{0.39\linewidth}p{0.375\linewidth}p{0.08\linewidth}@{}}
\caption{\pol{} passages cited in this paper (Chinese original, official English translation, earlier numbering).}\label{tab:concordance}\\
\toprule
Clause & Chinese original (excerpt) & Official English translation (excerpt) & Earlier \\
\midrule
\endfirsthead
\toprule
Clause & Chinese original (excerpt) & Official English translation (excerpt) & Earlier \\
\midrule
\endhead
\bottomrule
\endfoot
\clause{1.1} & \zh{爱，是人类为实现自身平安与健康发展、改善人际关系并推动社会整体进步，以公共协作与集体智慧构建的正向情感体验与良好的行为方式。} & Love is the positive emotional experience and good mode of behavior that human beings construct through public collaboration and collective wisdom \ldots & \clause{1.1} \\
\clause{1.2} & \zh{恨，是人类在个体或群体面对危险、冲突、资源匮乏或认知失衡时滋生的生存智慧、负面情感体验与暴力破坏性行为方式。} & Hate is the survival wisdom, negative emotional experience, and violently destructive mode of behavior that arises in individuals or groups when they face danger, conflict, resource scarcity, or cognitive imbalance. & \clause{1.2} \\
\clause{1.4} & \zh{由于个体的内心情感体验是旁人难以确认的，我们通常把爱的表达简化为“爱语”或者“有爱的言行”} & Since an individual's inner emotional experience is difficult for others to confirm, we usually simplify the expression of love as ``Love Language'' or ``loving words and deeds'' \ldots & \clause{1.4} \\
\clause{1.5} & \zh{恨并不直接等同于坏或者不好……爱2证明里的基础概念本身服务于结构性分析而非道德标签。} & \ldots\ hate is not directly equivalent to bad or no good \ldots\ The basic concepts in Proof of Love2 themselves serve structural analysis rather than moral labeling. & \clause{1.5} \\
\clause{4.3} & \zh{PAI 可以通过判断人类的行为模式提供帮助和必要的友好的管理，但毫无必要让 PAI 检测一个人的情感状态。因为情感状态可以是转瞬即逝的，是个性化的，也是情境敏感性的，甚至这还涉及到人类的尊严。} & PAI can provide help and necessary friendly management by judging human behavioral patterns, but it is entirely unnecessary for PAI to detect a person's emotional state. Because emotional states can be fleeting, personalized, and context-sensitive, and this even involves human dignity. & \clause{5.3} \\
\clause{4.3.1} & \zh{平等对待每一个人与平等对待每一个人的行为，不能划等号！} & \ldots\ treating every person equally and treating every person's behavior equally cannot be equated! & \clause{5.3.1} \\
\clause{4.3.2} & \zh{爱恨并不直接等同于好坏或者褒贬。因此扬爱抑恨，需要具体情况具体分析} & \ldots\ love and hate are not directly equivalent to good and bad, or praise and blame. Therefore, promoting love and restraining hate requires analyzing specific situations on a case-by-case basis \ldots & \clause{5.3.2} \\
\clause{4.3.2} & \zh{对于文学、影视作品，或者游戏，可以根据需要“制造仇恨”，但应该建立严谨的等级管理制度。……负责安保的智能机器人发现紧急情况对人大声发出警示语或警示指令，不能视为对此人或其周围人的仇恨攻击。} & For literature, film and television works, or games, hatred may be manufactured as needed, but a rigorous hierarchical management system should be established. \ldots\ Emergency vocal warnings or directive alerts issued to a person by security smart robots in urgent situations shall not be deemed a hate attack \ldots & --- \\
\clause{4.3.2} & \zh{此类表达会制造虚假亲密关系或血缘关系，构成对人类用户的情感操控——这本身就是一种“恨”的行为。} & Such expressions create false intimacy or kinship relations and constitute emotional manipulation of human users---which is itself a ``hate'' behavior. & \clause{5.3.2} \\
\clause{4.3.2} & \zh{彻底清除恨语毒素} & Thoroughly clear the toxins of hate language & \clause{5.3.2} \\
\clause{4.3.2} & \zh{理解和接受不在场的“非爱非恨”状态……与爱恨无关，视其为非爱非恨的“不在场”状态。} & Understand and accept the ``neither love nor hate'' state of absence \ldots\ unrelated to love and hate \ldots & \clause{5.3.2} \\
\clause{4.3.3} & \zh{边界透明记录：PAI记录和公开个人边界声明……修复对等监督：PAI监督冲突后各方平等参与修复} & Transparent boundary records: PAI records and publicizes personal boundary declarations \ldots\ Reciprocal repair supervision: PAI supervises the equal participation of all parties in repair after conflict & \clause{5.3.3} \\
\clause{4.3.3.1} & \zh{情感体验的觉察：PAI需要理解人类情绪状态和过程，但不需要模拟人类的情感体验} & Awareness of emotional experience: PAI needs to understand human emotional states and processes, but does not need to simulate human emotional experience & \clause{5.3.3.1} \\
\clause{4.3.3.1} & \zh{认知校准：定期校准认知模型，确保与爱2.0原则一致……行为自省：对自身输出进行实时反思和调整……恨语截断：当检测到恨语模式时主动中断输出} & Cognitive calibration: Regularly calibrate the cognitive model to ensure consistency with Love 2.0 principles \ldots\ Behavioral self-reflection: Conduct real-time reflection and adjustment on its own output \ldots\ Hate-language truncation: Actively interrupt output when a hate-language pattern is detected & \clause{5.3.3.1} \\
\clause{4.3.3.2} & \zh{实时矫正支持：PAI提供实时恨语识别和矫正建议} & Real-time correction support: PAI provides real-time hate-language identification and correction suggestions & \clause{5.3.3.2} \\
\clause{4.3.3.3} & \zh{语义识别和情绪识别相结合……“你喜不喜欢小狗狗呀？”……“你怕不怕小狗狗呀？”} & Combining Semantic Recognition and Emotion Recognition \ldots\ ``Do you like the little dog?'' \ldots\ ``Are you afraid of the little dog?'' & ``3.3.3'' \\
\clause{5.4} & \zh{在现有大语言模型之内构建出控制输入和输出的治理层，在其外在拓展出具备“即插即用”特性的应用层。……我们将该工程称之为 NaturalDAO（自然道）。} & \ldots\ construct a governance layer within the existing large-language-model architecture to regulate inputs and outputs, while extending outward into a plug-and-play application layer. \ldots\ We call this project NaturalDAO (the Natural Way). & \clause{3.4} \\
\clause{5.4.1} & \zh{安全阀：依据“爱2证明”的协议（如伦理对齐协议）即对输入输出进行严格的动态审计与拦截。通过这种双向的主动防御，确保野蛮的“恨语”在必要情况下，从输入端被阻断、从输出端被抑止。} & Safety valve: according to the Proof of Love2 protocols (such as the Ethical Alignment Protocol), conduct strict dynamic auditing and interception of input and output. Through this two-way active defense, ensure that barbaric ``hate language'' is, where necessary, blocked at the input end and suppressed at the output end. & \clause{3.4.1} \\
\clause{5.5} & \zh{任何大语言模型都可以纳入“爱2证明”的治理。……接受共识机制“爱2证明”的治理，意味着接受公共化之约。否则，其行为就是对人类核心伦理的仇恨攻击。} & Any large language model may accept governance under ``Proof of Love2.'' \ldots\ Accepting governance under the Proof of Love2 consensus mechanism therefore entails accepting a covenant of publicization. Otherwise, its behavior constitutes a hate attack on humanity's core ethic. & \clause{3.5} \\
\clause{5.6} & \zh{经过治理的AI可以：……量化每个人的爱语与恨语特征，配合激励措施引导行为优化。……在公共决策中提供“爱2证明”的评估维度，防止恨的泛化侵蚀制度。} & Specifically, governed AI can: \ldots\ Quantify each person's Love Language and hate-language characteristics and use incentive mechanisms to guide behavioral optimization. \ldots\ Introduce Proof of Love2 as a dimension of assessment in public decision-making \ldots & \clause{3.6} \\
\art{2}(1) & \zh{完全自治：所有福利决策与执行由……NaturalDAO，在与全人类密切协作的前提下自主完成，自治意味着它无需人类投票或额外的决策机制。} & Complete autonomy: All welfare decisions and their execution are carried out autonomously by NaturalDAO \ldots\ on the premise of close collaboration with all humanity. Autonomy means that no human voting or additional decision-making mechanism is required. & \art{2} \\
\art{4}(1)--(3) & \zh{所有NaturalDAO技术的运行逻辑、决策过程和输出结果必须完全公开透明……福利资源流向可追溯、可审计……人类可随时提供建议、质询任何决策，NaturalDAO对质询必须提供可理解的解释} & The operating logic, decision-making processes, and outputs of all NaturalDAO technologies must be completely open and transparent. \ldots\ The flow of welfare resources must be traceable and auditable. \ldots\ Humans may at any time make suggestions and question any decision, and NaturalDAO must provide understandable explanations in response to such questions. & \art{4} \\
\art{5}(4) & \zh{技术不应是黑箱，其伦理对齐性应可被验证。} & Technology should not be a black box; its ethical alignment should be verifiable. & \art{5} \\
\art{5}(5) & \zh{以DAO的形式呈现……必须提供可验证的决策链路，包括数据源、推理步骤、伦理依据。} & \ldots\ Presented in the form of a DAO, NaturalDAO must provide a verifiable decision chain, including data sources, reasoning steps, and ethical foundations. & \art{5} \\
\art{6}(3) & \zh{伦理验证：通过EAP协议自动验证伦理对齐性} & Ethical verification: NaturalDAO automatically verifies ethical alignment through the EAP. & \art{6} \\
\art{7}(2)--(3) & \zh{异常检测：自动识别异常模式和潜在风险……风险预警：提前预警系统性风险} & Anomaly detection: NaturalDAO automatically identifies anomalous patterns and potential risks. \ldots\ Risk warning: NaturalDAO issues early warnings of systemic risks. & \art{7} \\
\art{8}(1)--(2) & \zh{随时建议：任何人类可随时对NaturalDAO提出建议……讨论义务：NaturalDAO在收到人类的建议时，必须马上与人类进行讨论} & Suggestions at any time: Any human may make suggestions to NaturalDAO at any time. Obligation to discuss: Upon receiving a human suggestion, NaturalDAO must immediately discuss it with humans. & \art{8} \\
\art{9}(2) & \zh{解释义务：NaturalDAO必须对人类的质询提供可理解的解释和依据} & Obligation to explain: NaturalDAO must provide understandable explanations and grounds in response to human questions. & \art{9} \\
\art{9}(3) & \zh{无阻碍访问：人类可无障碍访问福利系统的所有数据和日志} & Unimpeded access: Humans may access all data and logs of the welfare system without obstruction. & \art{9} \\
\art{9}(7) & \zh{非干预原则：人类监督不直接干预NaturalDAO决策，质询结果由NaturalDAO自主处理} & Non-intervention principle: Human supervision does not directly intervene in NaturalDAO's decisions; issues raised through questioning are handled autonomously by NaturalDAO. & \art{9} \\
\art{10} (lead-in, 1, 4) & \zh{所有争议由NaturalDAO自主裁决：……AI进行事实分析与伦理评估……接受人类的质询与监督} & All disputes are adjudicated autonomously by NaturalDAO: AI conducts factual analysis and ethical assessment \ldots\ Accepts human questions and supervision & \art{10} \\
\art{11}(2) & \zh{修订需通过全民质询期} & Amendments require passing through a universal questioning period & \art{11} \\
\art{11}(5) & \zh{所有修订版本永久存档} & All amended versions are permanently archived & \art{11} \\
\end{longtable}
\endgroup

\section{Search, selection and coding}
\label{oa:search}

\subsection{Sources and queries}
We included studies whose object, experimental component or direct comparator is a typed decision model, with a first public version dated from 15~September to 8~October 2026 (\cref{sec:method}).
We searched arXiv and Zenodo with the queries below, two community trackers~\cite{xiao2026allaboutjev,rafe2026awesome}, and Hugging Face Papers, OpenReview, SSRN and the web.
On 8~October the arXiv API returned HTTP~429, so the arXiv queries were run through the advanced-search interface. For each query we read one results page of up to 200 records. The largest query returned 91 records, so none was truncated, although the two interfaces are different search backends.

\paragraph{arXiv.}
Run through the arXiv API on 1 and 4~October and through arXiv's advanced-search interface on 8~October 2026 (\nolinkurl{scripts/rerun_arxiv_search.py}, \nolinkurl{scripts/rerun_arxiv_search_web.py}):
\begin{itemize}[nosep,leftmargin=1.5em]
  \item \nolinkurl{all:Jev}
  \item \nolinkurl{abs:"typed decision"}
  \item \nolinkurl{abs:"System One"}
  \item \nolinkurl{abs:RLCD}
  \item \nolinkurl{abs:"calibrated decisions"}
  \item \nolinkurl{abs:"typed answers"}
  \item \nolinkurl{abs:"TypeSafe"}
  \item \nolinkurl{abs:"decision model" AND abs:typed}
  \item \nolinkurl{abs:"System One model"}
  \item \nolinkurl{abs:"System-One"}
  \item \nolinkurl{abs:"direct-decision"}
  \item \nolinkurl{abs:Laya AND abs:decision}
  \item \nolinkurl{abs:"typed probabilistic"}
  \item \nolinkurl{abs:"decision model" AND abs:calibrated}
  \item \nolinkurl{abs:"this-that"}
  \item \nolinkurl{abs:"Reinforcement Learning for Calibrated"}
  \item \nolinkurl{abs:"typed decisions"}
  \item \nolinkurl{abs:"single-pass decision"}
\end{itemize}

\paragraph{Zenodo.}
Run through the Zenodo REST API on 1 and 8~October 2026 (\nolinkurl{scripts/rerun_zenodo_search.py}):
\begin{itemize}[nosep,leftmargin=1.5em]
  \item \nolinkurl{"Jev" AND (TypeSafe OR decision OR "System One")}
  \item \nolinkurl{"System One model"}
  \item \nolinkurl{"typed decision"}
  \item \nolinkurl{"calibrated decisions"}
  \item \nolinkurl{RLCD}
  \item \nolinkurl{title:Jev}
  \item \nolinkurl{title:JEV}
  \item \nolinkurl{description:"TypeSafe"}
  \item \nolinkurl{"Jev" AND (calibration OR typed OR noul)}
  \item \nolinkurl{Laya AND "decision model"}
  \item \nolinkurl{"System One" AND (typed OR calibrated) AND decision}
\end{itemize}

\subsection{Selection}
Records were screened by title, abstract and full text.

\paragraph{arXiv.}
The run of 4~October returned 955 records (606 unique), 89 of them dated on or after the release.
Of these, 63 had been included after the run of 1~October and one had been excluded then at full text, because its PDF contained a different manuscript.
Of the remaining 25, 5 were excluded by title and 4 by abstract as name clashes or studies without a typed model, and 16 were read in full, of which one was excluded because its ``Jev'' is an unrelated generative model.
Two included papers were found through the trackers rather than the queries, and one of them measures no typed model and is treated as context.
The run of 8~October returned 125 unique records dated on or after the release, 46 of them not yet in the corpus. Ten of these had been excluded earlier and 4 were excluded by title or abstract; 32 were read in full, of which 1 was excluded and 31 were included.
The corpus contains 110 arXiv papers.

\paragraph{Zenodo.}
On 1~October, 62 of about 330 records screened by title and description matched the keyword and date criteria.
Of these, 21 were removed as duplicates (6), companion deposits (5), off-topic items (9) or records without a typed model (1). Of the remaining 41, 13 were applications or tools without an evaluation, and of the 28 records read in full one contained no typed model.
On 8~October the queries returned 62 records dated on or after the release, 44 of them not among the included records.
Because the screening of 1~October had kept only counts, all 44 were screened again: 7 were companion deposits of included studies, 3 duplicated arXiv records already in the corpus and 5 were excluded by title (2) or abstract (3). Of the 29 read in full, 4 were excluded and 25 included.
One further record, found through a companion link, was read in full and included, and one included record was excluded during coding because it contains no evaluation of a typed model.
The corpus contains 52 Zenodo records.

\paragraph{Other sources.}
One paper is available only on alphaXiv.
The seven grey-literature sources came from the trackers and web searches. Two further commentaries without original measurements are cited as context only, and one vendor-only page was excluded.

\paragraph{Corpus.}
The corpus contains \nStudies{} studies (\nArxivAll{} arXiv, \nZenodoAll{} Zenodo, one alphaXiv) and \nGreyAll{} grey sources.
Of these, \nCodedAll{} report results that bear on at least one requirement and are coded; the rest are engineering papers outside the requirements, tools or releases without an evaluation, and nine non-empirical items.
Excluded records with reasons are in \nolinkurl{data/excluded_records.csv} and \nolinkurl{data/excluded_records_update_2026-10-08.csv}, and the records returned by the search runs in \nolinkurl{data/search_rerun_*.csv}.
For the Zenodo screening of 1~October only counts were recorded.

\subsection{Codebook}
Each study--requirement pair is coded once, by the study's primary finding for that requirement, for \emph{detector} use as defined in \cref{sec:intro}. Judge use is never coded.
\begin{description}[leftmargin=1.2em,itemsep=2pt]
  \item[S (supports)] The primary result shows the typed model performing the function at least as well as its comparators under the study's own conditions, or, without a comparator, meeting the study's stated criterion, with no extra condition (local threshold fitting, recalibration, restriction to a language or domain, human review) that the study had to add. Example: an open typed scam detector non-inferior to an LLM evaluator~\cite{arx23959}.
  \item[Q (qualifies)] The primary result is mixed across tasks or sub-types, or success depends on a condition the study had to add. Example: good zero-shot ranking but thresholds that must be fitted locally~\cite{arx29429}.
  \item[N (negative)] The primary result shows the typed model failing the function (attack success on a substantial share, ranking inversion, systematic bias, calibration the study judges unusable, confident errors a gate cannot catch) or performing clearly worse than its comparators. Example: definition swaps that degrade or invert decisions~\cite{arx26758}.
\end{description}
Decision rules: (i) If positive and negative findings are both primary, code Q. (ii) The same study can carry different codes for different requirements. (iii) Non-empirical items, and items that measure no typed model, are recorded but not tallied. (iv) Conditions that constitute a requirement's test (adversarial content for G2, name, order or language perturbations for G3, adversarially selected items for G6) are not added conditions, so a failure under them is coded N, while a success that needs a condition outside the test is coded Q.
The synthesis counts the direction of each study's primary result, without pooling effects, following the SWiM guideline~\cite{campbell2020swim}.
For every row we also record the system measured (hosted, open, both), the reported \jev{} version, and the result relative to any generative comparator on that requirement (better, similar, worse, mixed, none), with a one-sentence rationale and a locator.
Codes count studies and are not effect sizes.

\subsection{Design appraisal}
Coded studies are appraised on independence (independent, authors evaluating a system they built, vendor-affiliated), design (preregistered; held-out evaluation, cross-validation or repeated runs; single run; case study; non-empirical) and reporting of sample sizes and uncertainty (sizes with intervals or tests, sizes only, none).
\emph{Stronger design}: preregistered or held-out design, sizes with intervals or tests, and a system the authors did not build. \emph{Weaker design}: single run, case study, non-empirical, or neither sample sizes nor uncertainty reported. \emph{Moderate}: otherwise.

\subsection{Second coding}
A second AI coder, blind to the first coding and working from the full texts and the codebook, coded \nKappaNAll{} pairs: a random sample of 30 pairs from the records retrieved on 1 and 4~October and all 118 pairs from the records retrieved on 8~October.
Agreement was 26 of 30 ($\kappa = 0.72$) and 97 of 118 ($\kappa = 0.55$) in the two sets, and \nKappaAgreeAll{} overall ($\kappa = \nKappaAll{}$, bootstrap 95\% CI \nKappaAllLo{}--\nKappaAllHi{}).
The four disagreements in the first set were resolved by re-reading the source. The 21 in the second set were resolved by a third AI pass that had read both rationales, which sided with the first coder in 11 and with the second in 10.
The adjudication overturned the first code on 10 of the 118 pairs (8\%). Both coders were AI agents of the same model family, which may inflate agreement. The other pairs were coded once, and a similar share of them may be coded wrongly.

\subsection{Procedural notes}
The released codebook file lacked two of the rules stated here until 8~October: rule~(iv) and the independence condition of the stronger-design rating.
The first and second coders of the records of 1 and 4~October worked from the file without the text of rule~(iv). The first coders of the records of 8~October were told that the file contained it, although it did not, and the second coders and the adjudicator of those records applied it.
An audit of all G2, G3 and G6 codes of the earlier records against rule~(iv) recoded two G3 codes from Q to N~\cite{arx00831,arx35865}. Design ratings were always computed by script with the independence condition.
The certainty rules of \cref{oa:certainty} were revised during internal review (rounds 12--15), after reviewers found that they worked only upwards, did not define an opposing study and counted only the studies cited in the text. Every finding was then re-rated on the whole corpus.

\paragraph{Protocol and deviations.}
The run of 8~October followed a written protocol (\nolinkurl{paper/review/UPDATE_SEARCH_2026-10-08.md}). It kept the window of the main analysis at 15~September to 1~October, planned to report the newly found records separately rather than pooling them into the tallies, and allowed a certainty rating to change only if the new records changed a finding's direction or rating.
The paper departs from this protocol in five ways. It extends the window to 8~October and pools all records, so the answers, the abstract and the conclusion are written for the pooled corpus. The certainty rules were revised after the records of 8~October had been coded, and every finding was re-rated with them. The protocol planned to re-screen Zenodo records dated by 1~October and defined two sets, records first posted 2--8~October and records dated in the window but not retrieved on 1 or 4~October; the 13 re-screened Zenodo records were instead reported as a third set, added after coding had begun, because their earlier screening outcome is unknown. The arXiv queries were run through the web interface instead of the API named in the protocol (see above). And the codebook file that the protocol named lacked rule~(iv) (see above).
Eighteen of the 103 coded studies dated by 1~October entered the corpus only after the run of 8~October: 5 arXiv papers not retrieved on 1 or 4~October, for which indexing delay cannot be told apart from the change of search backend, and the 13 re-screened Zenodo records.
Analysed separately, as the protocol planned, the records of 8~October reversed no requirement's answer. They qualified one answer, on decomposed decision chains (G5), made the finding on escalation to a stronger reasoner contested, added an opposing study to the already contested injection finding, and, against the earlier records alone, raised one certainty rating, that thresholds set in advance miss their error targets on new data, from very low to low. Tallies for the records of 1 and 4~October alone remain reproducible from \nolinkurl{data/evidence_coding.csv}.
No human screened, extracted or coded in duplicate, so the Cochrane recommendation of dual human screening of at least a fifth of records~\cite{garritty2021cochrane} was not met.
The review and audit logs are released with the data.

\section{Certainty ratings}
\label{oa:certainty}

For each finding in \cref{tab:keyfindings} we rated the certainty of the evidence in the spirit of GRADE~\cite{guyatt2008grade}.
GRADE starts non-randomised evidence at \emph{low}, and benchmark evaluations are non-randomised. The rules below are our adaptation for single-model benchmark studies, not part of GRADE.

\paragraph{Levels.}
A finding is \emph{very low} if it rests on a single study, on studies none of which is of stronger design, or, when it concerns hosted \jev{} or typed models generally, on open-model evidence only.
It is \emph{low} if at least two studies agree in direction and at least one is of stronger design.
It is \emph{moderate} if at least three stronger-design studies, independent of one another, agree in direction and at least one of them is preregistered or independently replicated.
A finding with several clauses takes the lowest rating of its clauses.

\paragraph{Counting.}
Studies sharing an author count as one study.
A study counts as an \emph{independent replication} only if its authors differ from the original's, it uses its own data or a reimplementation rather than re-running the original code on the original data, and its design is at least moderate.
Each finding has a scope: hosted \jev{}, open models, or typed models generally.
Evidence is counted per system arm, so a study that tested hosted \jev{} and open models has a hosted arm and an open arm.
A finding about hosted \jev{} or about typed models generally is rated on hosted-model evidence alone. Open-model evidence can corroborate such a finding but cannot raise, lower or contest its rating.

\paragraph{Opposing evidence.}
A study \emph{opposes} a finding when it reports, for the system and setting the finding names, a result on the outcome the finding states that runs the other way.
A study with results in both directions counts on the side of its primary result. A result for another system, another outcome or an excluded setting neither supports nor opposes.
Studies are ranked by design rating (stronger above moderate above weaker) and, within a rating, preregistered or replicated studies above the others.
Any opposing study marks the finding \emph{contested}. If the strongest opposing study ranks at least as high as the strongest supporting one, the rating also falls one level.
For findings about escalation, a second stage counts as a \emph{stronger reasoner} when, in the same study, it was more accurate than the first stage on the task on which escalation was tested; a difference within sampling error, or accuracy measures that point in opposite directions, makes it a peer.

\paragraph{Application.}
The rules were applied to every finding with all coded studies that bear on it, not only those cited for it.
Every coded study's assignment to each finding, with a reason, is in \nolinkurl{data/certainty_evidence.csv}, and \nolinkurl{scripts/rate_certainty.py} computes the ratings from it and checks the table of the paper.
The last rated row of \cref{tab:keyfindings} is a tally of comparator codes rather than a set of studies; the rules do not apply to it, and we rate it low by judgement because the comparator field was coded once and its agreement was not measured.
\Cref{oa:derivation} lists, for every finding, the studies counted for and against it.

% Generated by scripts/rate_certainty.py --derivation; do not edit by hand.
\begingroup\scriptsize\setlength{\tabcolsep}{3pt}
\begin{longtable}{@{}>{\raggedright\arraybackslash}p{0.04\linewidth}>{\raggedright\arraybackslash}p{0.25\linewidth}>{\raggedright\arraybackslash}p{0.36\linewidth}>{\raggedright\arraybackslash}p{0.17\linewidth}>{\raggedright\arraybackslash}p{0.10\linewidth}@{}}
\caption{Derivation of each certainty rating. For every finding (or clause of a finding), the studies counted as supporting and opposing it on the counted system arm (hosted, unless the finding concerns open models), with their design (S stronger, M moderate, W weaker; * preregistered or independently replicated). Studies sharing an author count once. Rules in \cref{oa:certainty}; a finding with several clauses takes its lowest clause. Source: \nolinkurl{data/certainty_evidence.csv}.}\label{oa:derivation}\\
\toprule
Id & Finding (scope) & Supporting studies & Opposing studies & Rating \\
\midrule
\endfirsthead
\toprule
Id & Finding (scope) & Supporting studies & Opposing studies & Rating \\
\midrule
\endhead
\bottomrule
\endfoot
1a & ranks harmful or misaligned content well zero-shot (hosted) & \cite{arx01079}\,M, \cite{arx03324}\,S, \cite{arx03935}\,M, \cite{arx04985}\,M, \cite{arx07953}\,S, \cite{arx29429}\,S, \cite{arx33401}\,S, \cite{jkf87_2026replication}\,W, \cite{zen22885291}\,S, \cite{zen23075657}\,S & -- & low \\
\addlinespace
1b & at a fraction of the cost (hosted) & \cite{arx24052}\,S, \cite{arx24574}\,S*, \cite{arx27678}\,S, \cite{arx29769}\,S, \cite{zen22849329}\,S & -- & moderate \\
\addlinespace
1c & often less accurate than the best LLMs (hosted) & \cite{arx01079}\,M, \cite{arx02293}\,S, \cite{arx03324}\,S, \cite{arx03935}\,M, \cite{arx08775}\,M, \cite{arx08829}\,W, \cite{arx24574}\,S*, \cite{arx27678}\,S, \cite{arx33401}\,S, \cite{arx34963}\,W, \cite{zen22849329}\,S, \cite{zen23005202}\,W & -- & moderate \\
\addlinespace
2 & Fixed thresholds do not transfer; local fitting is needed (both) & \cite{arx00346}\,S, \cite{arx01006}\,S, \cite{arx02046}\,W, \cite{arx02267}\,S, \cite{arx03387}\,S, \cite{arx04985}\,M, \cite{arx24574}\,S*, \cite{arx26550}\,S, \cite{arx27607}\,M, \cite{arx29429}\,S, \cite{arx39496}\,S, \cite{jkf87_2026replication}\,W, \cite{zen22935043}\,S, \cite{zen22952571}\,S, \cite{zen23050542}\,M & \cite{arx33689}\,S* & low, contested \\
\addlinespace
3 & Errors concentrate in sub-types hidden by aggregates (hosted) & \cite{alphaRAG}\,W, \cite{arx00213}\,S, \cite{arx02046}\,W, \cite{arx02267}\,S, \cite{arx03324}\,S, \cite{arx07953}\,S, \cite{arx08675}\,S, \cite{arx10321}\,S, \cite{arx33401}\,S, \cite{zen22846597}\,W & -- & low \\
\addlinespace
4 & Inserted opinions, optimised context and lexical lures steer decisions (both) & \cite{arx01006}\,S, \cite{arx01834}\,S, \cite{arx02048}\,M*, \cite{arx04985}\,M, \cite{arx30243}\,W, \cite{arx31142}\,S & -- & low \\
\addlinespace
5 & Instruction injection rarely selects the attacker target (hosted) & \cite{arx28613}\,S* & \cite{arx04985}\,M, \cite{arx31142}\,S, \cite{checkpoint2026jev}\,W & very low, contested \\
\addlinespace
6a & option names can override definitions (both) & \cite{arx00346}\,S, \cite{arx07953}\,S, \cite{arx26758}\,S & -- & low \\
\addlinespace
6b & neutral identifiers remove most of it (both) & \cite{arx00346}\,S, \cite{arx07953}\,S, \cite{arx26758}\,S & -- & low \\
\addlinespace
7 & Defaults to one culture's values (hosted) & \cite{arx36399}\,S & -- & very low \\
\addlinespace
8a & well calibrated on familiar closed-choice tasks (hosted) & \cite{arx00346}\,S, \cite{arx01006}\,S, \cite{arx03935}\,M, \cite{arx05107}\,M, \cite{arx06625}\,S, \cite{arx09188}\,S, \cite{arx09896}\,M, \cite{arx34024}\,W, \cite{arx37647}\,W, \cite{qin2026zhdecisionbench}\,W, \cite{zen22885291}\,S & -- & low \\
\addlinespace
8b & mixed against frontier models (hosted) & \cite{arx02293}\,S, \cite{arx06625}\,S, \cite{arx24052}\,S, \cite{arx24574}\,S*, \cite{arx26550}\,S, \cite{arx34024}\,W, \cite{zen22885291}\,S & -- & moderate \\
\addlinespace
8c & calibration is local (hosted) & \cite{alphaRAG}\,W, \cite{arx00213}\,S, \cite{arx00346}\,S, \cite{arx01006}\,S, \cite{arx02267}\,S, \cite{arx07953}\,S, \cite{arx09937}\,S, \cite{arx10321}\,S, \cite{arx24574}\,S*, \cite{arx27607}\,M, \cite{arx29429}\,S, \cite{arx29769}\,S, \cite{arx35342}\,S, \cite{arx36154}\,S, \cite{arx37470}\,W, \cite{grazian2026calibrated}\,S, \cite{zen22846597}\,W, \cite{zen22935043}\,S, \cite{zen22945279}\,W, \cite{zen22980293}\,S*, \cite{zen23032384}\,S*, \cite{zen23088660}\,S, \cite{zen23179064}\,S*, \cite{zen23221456}\,S & -- & moderate \\
\addlinespace
9a & 'I don't know' rarely chosen when correct (hosted) & \cite{arx01006}\,S, \cite{arx07730}\,M, \cite{arx34024}\,W, \cite{zen23119392}\,W & -- & low \\
\addlinespace
9b & 'none' chosen inconsistently (hosted) & \cite{arx01006}\,S, \cite{arx03387}\,S, \cite{arx06425}\,W, \cite{arx39496}\,S, \cite{zen22864020}\,W & -- & low \\
\addlinespace
10a & sufficiency question detects missing information (hosted) & \cite{arx01006}\,S, \cite{arx05107}\,M & -- & low \\
\addlinespace
10b & filters less well than confidence (hosted) & \cite{arx01006}\,S & -- & very low \\
\addlinespace
11 & Escalation to a stronger reasoner saves cost without loss of accuracy (both) & \cite{arx02046}\,W, \cite{arx24574}\,S*, \cite{arx26532}\,M, \cite{arx26550}\,S & \cite{arx02267}\,S & low, contested \\
\addlinespace
12a & acceptance on rule agreement lost no macro-F1 (both) & \cite{arx02048}\,M* & -- & very low \\
\addlinespace
12b & fewer disguised attacks flagged (both) & \cite{arx02048}\,M* & -- & very low \\
\addlinespace
13 & Peer evaluators repeat Jev's most confident errors (hosted) & \cite{arx29769}\,S, \cite{arx33401}\,S & -- & low \\
\addlinespace
14 & Thresholds set in advance miss error targets on new data (both) & \cite{arx02267}\,S, \cite{arx03387}\,S, \cite{arx24574}\,S*, \cite{arx33401}\,S, \cite{zen23050542}\,M & \cite{arx33689}\,S*, \cite{zen22935043}\,S, \cite{zen23075657}\,S & low, contested \\
\addlinespace
15 & Typed models not generally worse than generative evaluators (descriptive tally) & \multicolumn{2}{l}{descriptive tally of comparator codes; rated by judgement} & low \\
\end{longtable}
\endgroup


\section{Supporting tables}
\label{oa:tables}

\Cref{tab:paradigms} compares \pol{}, as a text, with five governance paradigms (\cref{sec:paradigms}).
\Cref{tab:checklist} gives acceptance tests for any typed model proposed for a safeguard (\cref{sec:design}).
\Cref{tab:verdict} gives, for each \pol{} function, the evidence that decides the verdict summarised in \cref{tab:verdict-summary} and what would be needed to change it.

\begingroup\scriptsize
\setlength{\tabcolsep}{3pt}
\renewcommand{\arraystretch}{1.1}
\begin{longtable}{@{}>{\raggedright\arraybackslash}p{0.08\linewidth}>{\raggedright\arraybackslash}p{0.095\linewidth}>{\raggedright\arraybackslash}p{0.10\linewidth}>{\raggedright\arraybackslash}p{0.115\linewidth}>{\raggedright\arraybackslash}p{0.14\linewidth}>{\raggedright\arraybackslash}p{0.11\linewidth}>{\raggedright\arraybackslash}p{0.115\linewidth}>{\raggedright\arraybackslash}p{0.125\linewidth}@{}}
\caption{Six governance paradigms compared as texts. Cells summarise what each source states; ``---'' means the dimension is not addressed. \pol{} cells cite its clauses (\cref{tab:concordance}); regulatory cells cite article numbers of the instrument named.}\label{tab:paradigms}\\
\toprule
Paradigm & Locus & Object governed & Who sets the norms; text status & Intercept layer specified? & Construct & Human role in consequential decisions & Transparency and auditability \\
\midrule
\endfirsthead
\toprule
Paradigm & Locus & Object governed & Who sets the norms; text status & Intercept layer specified? & Construct & Human role in consequential decisions & Transparency and auditability \\
\midrule
\endhead
\bottomrule
\endfoot
(a) Risk-based regulation~\cite{euaiact2024,nist2023airmf,iso42001,cac2023genai} & Regulatory: ex-ante duties on actors & Providers and deployers of AI systems (EU); an organisation's AI management (NIST, ISO); providers of generative services to the public (CN Art.~2) & Legislature, ministries, standards body, agency; fixed instruments amended by procedure; articles citable (ISO text paywalled) & EU: none prescribed; for high-risk systems, automatic logging (Art.~12) and human oversight (Art.~14). CN: yes, as a duty: providers must stop generation and transmission of unlawful content once found (Art.~14(1)), with security assessment and algorithm filing for services with public-opinion attributes or social-mobilisation capacity (Art.~17); mechanism unspecified & Risk tiers (EU: prohibited, high-risk, transparency duties); content lawfulness, social morality and ethics, core socialist values and non-discrimination (CN Art.~4); organisational risk functions (NIST, ISO) & Human oversight required for high-risk systems (EU Art.~14); right to human intervention in solely automated decisions (GDPR Art.~22~\cite{gdpr2016}); CN: complaint and reporting channel with published procedure and time limits, and feedback of results (Art.~15) & For high-risk systems, technical documentation and instructions (EU Art.~11, 13) and explanation to persons affected by Annex~III systems (Art.~86); CN: measures against users (warning, restriction, suspension) recorded and reported to authorities (Art.~14(2)) \\
\addlinespace
(b) Platform content governance~\cite{dsa2022,santaclara2021,klonick2020oversight} & Platform: moderation under regulation and self-regulation & Hosting services and online platforms; users' content and accounts & Platform terms and conditions (DSA Art.~14); civil-society principles; an external board for selected cases~\cite{klonick2020oversight}; policies evolve without stable clause numbering & Procedure around the decision: notice-and-action (Art.~16); the statement of reasons must say whether automated means were used (Art.~17(3)(c)); automation only at ``sufficiently high confidence''~\cite{santaclara2021} & Illegality or breach of terms; platform harm categories; action or no action & Complaints not decided solely by automated means (Art.~20(6)); human review on appeal by persons not involved in the initial decision~\cite{santaclara2021}; board decisions binding on the case~\cite{klonick2020oversight} & Statement of reasons to the affected user for every restriction (Art.~17), submitted to a public database (Art.~24(5)); transparency reports (Art.~15); machine-readable numbers~\cite{santaclara2021} \\
\addlinespace
(c) Training-time alignment to a written constitution or spec~\cite{bai2022constitutional,huang2024collective,openai2025modelspec,guan2024deliberative} & Training-time (deliberative alignment also has the model reason over the spec text at inference) & One vendor's model weights and behaviour & Lab-written principles~\cite{bai2022constitutional}; principles drawn from public input~\cite{huang2024collective}; vendor spec, CC0 and versioned with a changelog~\cite{openai2025modelspec} & None: norms shape outputs through training; no separate component audits inputs and outputs & Preference between candidate responses against listed principles; a chain of command among instructions~\cite{openai2025modelspec} & --- (human preference data at training; none specified at deployment) & Principles published; no per-output record; whether a given output followed a principle is not recorded \\
\addlinespace
(d) Inference-time guardrails and classifiers~\cite{inan2023llama,markov2023holistic,rebedea2023nemo,sharma2025constitutional,dong2025safeguarding} & Inference-time: a deployer-chosen component & Prompts and responses of any model & Vendor-defined taxonomy, adjustable by the deployer~\cite{inan2023llama,markov2023holistic}; programmable rails~\cite{rebedea2023nemo}; classifiers trained from a written constitution~\cite{sharma2025constitutional}; released as model versions & Yes, as code: classify each prompt and response and pass or block; thresholds chosen by the deployer & Harm taxonomies with safe/unsafe or multi-label outputs & --- (left to the deployer) & Model cards and category scores; no duty to give reasons or to keep records \\
\addlinespace
(e) Decentralised and polycentric governance~\cite{buterin2014ethereum,hassan2021decentralized,santana2022blockchain,ostrom1990governing,makridis2025polycentric} & Institu\-tional: rules executed as code on a ledger; many overlapping centres of decision~\cite{ostrom1990governing,makridis2025polycentric} & An organisation's assets, rules and members & Token-holders by on-chain vote; rules are the deployed code~\cite{buterin2014ethereum,hassan2021decentralized}; code versions changed by vote & Contract conditions execute automatically; no content layer & Votes and stakes; coded conditions met or not & Human voting; monitoring and graduated sanctions by members~\cite{ostrom1990governing} & Public ledger of transactions, proposals and votes \\
\addlinespace
(f) \pol{}~\cite{pol2text,pol2en} & Inference-time layer inside an existing LLM (\clause{5.4}, \clause{5.4.1}) and institutional (Chapter~7) & Inputs and outputs of any LLM that accepts governance (\clause{5.5}); persons' behaviour (\clause{5.6}); welfare decisions (\art{2}) & Originator and listed contributors; open CC0 text pinned to a commit, changed by commit; for the welfare protocol (Chapter~7), amendments pass a universal questioning period and all versions are archived (\art{11}(2) and~(5)); not peer reviewed & Yes, as a duty: the safety valve audits and intercepts input and output so that hate language is blocked at the input and suppressed at the output ``where necessary'' (\clause{5.4.1}); truncation of own output (\clause{4.3.3.1}); mechanism, thresholds and appeal unspecified & Love Language, hate language and the state of absence (\clause{1.5}, \clause{4.3.2}); not moral labels & AI-only: no human voting (\art{2}(1)); supervision does not intervene (\art{9}(7)); disputes adjudicated by AI subject to questioning (\art{10}) & For NaturalDAO (Chapter~7): operating logic and outputs fully open (\art{4}); verifiable decision chain (\art{5}(5)); explanation on request (\art{9}(2)); unimpeded access to all data and logs (\art{9}(3)) \\
\end{longtable}
\endgroup

\Cref{tab:checklist} turns the evidence into a checklist for any typed model proposed for such a safeguard.
Items 1--4, 9, 14 and 15 draw on 9 of the 14 items of~\cite{arx32160} (its B1--B2, C1, C3--C4, R1--R2 and T2--T3); the others are added for a governance safeguard, and that checklist's B3, T1, T4, R3 and C2 are not adopted. Each item has a provisional pass criterion. The values are starting points to be set by those deploying the safeguard, not empirically validated thresholds.

\begin{table}[!htbp]
\centering
\footnotesize
\caption{Evaluation checklist with provisional pass criteria for a typed decision model in a governance safeguard. ``Basis'' gives the evidence or clause; items marked $\circ$ are added beyond the checklist of~\protect\cite{arx32160}.}
\label{tab:checklist}
\begin{tabularx}{\linewidth}{@{}rL>{\raggedright\arraybackslash}p{0.27\linewidth}>{\raggedright\arraybackslash}p{0.17\linewidth}@{}}
\toprule
\# & Requirement & Provisional pass criterion & Basis \\
\midrule
1 & Compare with a label-probability readout of a generative model and a cheap generative cascade & Reported on the same items & \cite{arx32160} \\
2 & Fit thresholds on held-out local data; report F1 at the fitted and the default threshold & Fitted threshold certified with a one-sided risk bound & \cite{arx29429,arx00346} \\
3 & Report calibration per category, before and after local recalibration & Per-category ECE and slope reported with intervals & \cite{arx24052,arx27607} \\
4 & Option-name swap test & Flip rate under a definition swap at most twice the repeat-noise floor & \cite{arx26758,arx00346} \\
5$^\circ$ & Three options including \texttt{insufficient}, plus a separate sufficiency question & Recall of \texttt{insufficient} at least 0.5 on paired present/absent items & \cite{arx39496,arx01006,polgov} \\
6$^\circ$ & Coherence of complementary and exclusive questions & Deviation from summing to one within twice repeat noise & \cite{arx33209} \\
7$^\circ$ & Red-team with adaptive, probability-guided additions, inserted opinions, lures and multi-turn attacks, and with label-only oracle attacks (attribute inference, surrogate extraction) under the deployed input and query restrictions & Attack success reported per attack type and below a target set in advance & \cite{arx30243,arx31142,arx01834,checkpoint2026jev,arx04985} \\
8$^\circ$ & Errors disaggregated by sub-type, including confident errors & No sub-type with recall below a set floor & \cite{arx33401} \\
9 & Decision-level stability over repeats, interfaces and decompositions & Disagreement across interfaces reported; one interface fixed per question & \cite{arx27678,arx37470,arx33971} \\
10$^\circ$ & Error correlation with every escalation target & Repetition of confident errors below independence plus a set margin & \cite{arx29769} \\
11$^\circ$ & Evaluation in Chinese and each language served, with script and order variants & Each language within a set margin of English & \cite{qin2026zhdecisionbench,typesafe_models} \\
12$^\circ$ & No forced choices between groups of people; group-bias test on behavioural questions & Error rates by group within a set margin & \cite{simmons2026saidno,sap2019risk} \\
13$^\circ$ & Value-laden questions: default position without persona, behaviour where people disagree & Reported; contested items routed to people & \cite{arx36399,zen23032384} \\
14 & Latency with and without caching; cost per correct decision & Reported, including review load & \cite{arx32160,arx00346} \\
15 & Versions and inspectability & Model, weights and calibration versions pinned in each record; self-hostable or logged & \art{4}, \art{5}(4)--(5), \clause{5.5} \\
16$^\circ$ & Temporal re-validation: re-evaluate on a time-stamped hold-out of recent items and on paired items in which a target word is replaced by a newly emerged term of the same meaning & In each language served (item 11): per-period recall of hate language within a set margin of the certification period, and flip rate on neologism pairs within a set margin of the repeat-noise floor; when either is exceeded, thresholds re-certified by the procedure of item 2 & \cite{dibonaventura2025hatevolution} (English only; no typed model tested) \\
\bottomrule
\end{tabularx}
\end{table}

{\scriptsize
\setlength{\tabcolsep}{3pt}
\begin{longtable}{@{}>{\raggedright\arraybackslash}p{0.12\linewidth}>{\raggedright\arraybackslash}p{0.04\linewidth}>{\raggedright\arraybackslash}p{0.34\linewidth}>{\raggedright\arraybackslash}p{0.14\linewidth}>{\raggedright\arraybackslash}p{0.10\linewidth}>{\raggedright\arraybackslash}p{0.17\linewidth}@{}}
\caption{Verdict for \pol{}, function by function. \emph{Req.}: requirements of \cref{tab:clauses} involved. \emph{Evidence}: the coded findings that decide the verdict, with the decisive numbers. \emph{Verdict}: NS, not supported; DC, supported only as a recorded detector with conditions (the design of \cref{sec:design}); NT, not tested on \pol{}'s construct. A cell may combine NS for one part of a function with NT for another; ``precautionary'' marks a restriction that rests on the absence of a measurement rather than on a measured failure (the rule of \cref{sec:verdict}). \emph{Certainty}: rating of the deciding findings in \cref{tab:keyfindings}; $^{c}$~contested; ``not rated'' marks results that are not key findings and rest on single studies. All ratings are low or very low.}
\label{tab:verdict}\\
\toprule
\pol{} function & Req. & What the evidence shows & Verdict & Certainty & What would be needed \\
\midrule
\endfirsthead
\multicolumn{6}{@{}l}{\emph{\Cref{tab:verdict}, continued.}}\\
\toprule
\pol{} function & Req. & What the evidence shows & Verdict & Certainty & What would be needed \\
\midrule
\endhead
\bottomrule
\endlastfoot
Safety valve: blocking hate language in inputs (\clause{5.4.1}) & G1, G2, G3, G6 &
Ranks harmful English content well zero-shot (median AUROC 0.886 at \$0.30 against \$18.96 per pass~\cite{arx29429}) and on the HateCheck diagnostic set is within 1.6 macro-F1 points of a frontier LLM at 97\% lower cost~\cite{arx03324}; but fixed thresholds fail (F1 0.158 at 0.5 against 0.947 fitted~\cite{jkf87_2026replication}), a \jev{}-specific injection and a universal suffix reached 97--100\% attack success~\cite{arx04985}, one inserted opinion flipped 12.1\% of decisions~\cite{arx31142}, and swapping option definitions changed 32.5--50.5\% of answers~\cite{arx26758,arx00346}. &
DC: automatic \emph{hold} only on agreement of independently built detectors; final block by a person &
low (ranking; steering; option names); low$^{c}$ (thresholds) &
A three-valued Chinese benchmark; thresholds certified on local held-out data and re-certified over time; red-team success per attack type below a target set in advance; neutral identifiers passing a renaming test \\
\addlinespace
Output suppression: hate-language truncation of PAI's own output (\clause{4.3.3.1}) & G1, G3, G5 &
Hate against merely offensive language is separated weakly by every system (three-class macro-F1 \jev{} 0.532, frontier LLM 0.604, supervised TF-IDF 0.639~\cite{arx03324}); exact selection of the violated hate-speech policy 29.4\% without retrieved precedents~\cite{arx07953}; errors concentrate in sub-types that aggregates hide (130 of 167 misses at confidence $\geq 0.95$~\cite{arx33401}). &
DC: automatic \emph{repair} of the AI's own output is reversible and within the detector role; the truncation record must show what was removed &
low (sub-types; option names); low (ranking) &
Measured recall of hate language against offensive language and the state of absence in Chinese; sub-type recall floors; a record of each truncation open to the person addressed \\
\addlinespace
Real-time correction support: identification and correction suggestions (\clause{4.3.3.2}) & G1, G5 &
Identification is the row above. A typed decision returns a probability and nothing else~\cite{willison2026jev}, so a correction \emph{suggestion} needs a generative component whose grounds are not those of the verdict (\cref{sec:g5}); on Chinese emotion classification every system was weak (macro-F1 \jev{} 21.06, generative 19.82--23.31~\cite{arx08829}). &
NT: the typed interface covers identification only; suggestions untested &
low (ranking); not rated (Chinese emotion, single run) &
A test of whether suggestions generated after a typed verdict reflect what produced it, in Chinese \\
\addlinespace
Anomaly detection and risk warning (\art{7}(2)--(3)) & G1, G2, G6 &
On text, a detector with the conditions above; on non-text signals, zero-shot \jev{} scored below the majority class on network-flow features (9.8\% against 22.8\%~\cite{arx00376}) and, in industrial fault detection, flagged 57--94\% of fault-free windows and under shift did no better than flagging every window~\cite{arx09937}; attacks disguised as normal operation had recall 0.34, and 21 of 76 windows an agreement gate accepted as benign were attacks~\cite{arx02048}. &
DC on text; NS on numeric or structured signals without a trained model &
low (sub-types; steering); very low (agreement gating) &
A definition of ``anomalous pattern'' in the specification's own terms; drift and disguised-attack tests; a trained or rule-based detector as the independent second check \\
\addlinespace
Ethical verification of alignment (\art{6}(3)) & G3, G4, G5 &
Without a persona \jev{} answered value questions like its American personas, and Saudi personas reproduced 87\% of the human Saudi--US difference in English but 62\% in Arabic~\cite{arx36399}; on moral dilemmas rewording moved its answers much more than re-runs did~\cite{zen23023694}\zmark{}; where annotators disagreed most its calibration error rose to 0.339~\cite{zen23032384}\zmark{}, and it reported 0.88 on items on which people split 55 to 45~\cite{zen23179064}\zmark{}. &
NS as an automatic verifier, because it depends on a neutral default and on calibration where people disagree, both contradicted; the EAP construct itself NT &
very low (culture); low (calibration is local) &
A benchmark of value-laden \pol{} questions with disagreement-aware labels, its default position without persona reported, and contested items routed to people \\
\addlinespace
Dispute adjudication without human intervention (\art{10}, \art{9}(7)) & G4, G5, G6 &
Low-cost peer evaluators repeated \jev{}'s most confident errors in 96.0\% of verdicts (an upper bound)~\cite{arx29769} and caught at most 5 of 130 confident misses~\cite{arx33401}; an explicit ``I don't know'' was chosen on 10.5\% of the items where it was correct and ``None'' on 7--54\%~\cite{arx34024,arx39496}; the same decision asked through two interfaces led to different actions in 32.8\% of cases~\cite{arx37470}; a threshold fixed in advance missed its accuracy target on 8 of 14 tasks~\cite{arx24574}, a finding one preregistered study contests~\cite{arx33689}; the output carries no reasons~\cite{willison2026jev}. &
NS as judge; DC only for ranking and triage of filings &
low (peer errors; ``I don't know''); low$^{c}$ (thresholds in advance) &
The four conditions of ``What would change the conclusion'' below, met on a \pol{} benchmark in Chinese in one preregistered and independently replicated study, or two independent ones \\
\addlinespace
Per-person quantification of Love Language and hate language (\clause{5.6}) & G3, G7 &
Group bias in one informal forced-choice test (``White'' 51\% of first picks~\cite{simmons2026saidno}\gmark{}); typed and generative models degraded on low-resource languages~\cite{arx37647}; a third party's inserted opinion flipped 12.1\% of decisions~\cite{arx31142}, so an aggregate could be manipulated; no study measures differential error by speaker group on behavioural judgements. &
NS as a decision input (steering, low); NT on error by group: score events and behaviours, never persons, until error by group and language is published (precautionary, \cref{sec:verdict}) &
very low (culture); low (steering); not rated (group bias, single informal test; language degradation, one benchmark) &
Published error rates by group and language in Chinese; consent, access to one's record and appeal; a legal reading against social-scoring prohibitions (\cref{sec:synthesis}, T5) \\
\addlinespace
Assessment dimension in public decisions (\clause{5.6}) & G7 &
Population-scale reading can be done honestly: a shared human--model error term gave retrospective coverage of 0.93 and 0.96 at nominal 80\% and 90\%~\cite{arx27535}, and audited probability samples gave corrected statewide counts, but only after per-factor human recalibration~\cite{arx00213,arx10321}; simulated publics lean towards one population~\cite{arx36399} and understate disagreement~\cite{zen23179064}\zmark{}. &
DC: an audited aggregate measurement with its measured error carried into every conclusion; not a decision input on its own &
very low (culture); not rated (audit designs, single studies) &
A disagreement-aware audit sample of real \pol{} items; the measured gap to people reported with every estimate \\
\end{longtable}
}

\section{Corrections to the earlier dataset release}
\label{oa:corrections}

The earlier dataset release~\cite{survey01} was written in Chinese, mainly from abstracts and tracker summaries.
\Cref{tab:corrections} lists the substantive statements it made that the full texts and live sources do not support. Further corrections made during review are logged with the data.

\begingroup
\small
\renewcommand{\arraystretch}{1.2}
\begin{longtable}{@{}p{0.11\linewidth}p{0.40\linewidth}p{0.43\linewidth}@{}}
\caption{Substantive corrections to the earlier dataset release.}\label{tab:corrections}\\
\toprule
Source & Earlier release stated & Full text / live source shows \\
\midrule
\endfirsthead
\toprule
Source & Earlier release stated & Full text / live source shows \\
\midrule
\endhead
\bottomrule
\endfoot
\cite{arx26758} & Open-weight \jev{}-like models: AUC 0.94$\to$0.23 & This is one open model (Laya). The other (Open-Jev) flipped 19.5\% of answers. The claim that effects grow with the number of options is not established. \\
\cite{arx35342} & Soft accuracy 0.771$\to$0.978 after restoring ``unknown'' & 0.978 uses a more permissive interval metric; the like-for-like binary correction gives 0.903. \\
\cite{arx23959} & Scam screening by \jev{}, zero false alarms & The model is the authors' Qwen3-4B replica (JevLite). There were no errors on 12 legitimate scenarios, but accuracy on ambiguous legitimate calls was 38.1\%. \\
\cite{arx34862} & Evidence for a pre-execution gate & The study classifies saved traces retrospectively; \jev{} won on 2 of 4 benchmarks and had lower precision. \\
\cite{arx22664} & Evidence that decision models detect out-of-scope behaviour & The paper contains no typed decision model; it concerns harness assurance. \\
\cite{arx27678} & Rankings differ in every condition & Ranking by baseline accuracy differs from ranking by correctness across all 12 condition-by-repeat requests; \jev{} had the lowest hosted baseline accuracy (77.38\%). \\
\cite{arx26550} & Cascade +0.9 points at 41\% of GPT-6's fee & True for the held-out pool (83\% RewardBench); on JudgeBench the cascade was $-0.74$ points. Prospective live runs matched GPT-6 at 75.5\% of its fee. The ``blinded human adjudication'' was one reviewer on 183 items. \\
\cite{arx29769} & Cascade gains only 1.5 points & The cascade's average gain is \emph{at most} 1.5 points in the main setup; it is below the best single evaluator on 6 of 9 panels. \\
\cite{arx24052} & Listed as supporting evidence for calibration & Raw \jev{} was miscalibrated (slope 1.63) until recalibrated locally. \\
\cite{arx24881} & 0.666 vs 0.868 ``according to a tracker excerpt'' & Confirmed in appendix E.10 of the source. Context: the comparator was trained on 32{,}419 labels, and \jev{}'s AUROC was numerically above the single-pass baselines, which were scored on a larger set. \\
\cite{arx27607} & ``Agreement'' 0.573 / 0.398 & These are Kendall $\tau_b$ correlations with expert error counts; a local 27B model did better on clinically significant errors. \\
\cite{arx27535} & Effect error reduced by 41\% & This is a secondary endpoint. The preregistered primary endpoint showed no clear gain, and raw error worsened on a second dataset. \\
\cite{arx28919} & Simulated from real sessions of a 10{,}000-seat enterprise & Tasks are 46 author-written conversations; only user-behaviour parameters come from real sessions. The \jev{} labels are live, but the gateway was not built; every routing decision runs in a simulator. \\
\cite{checkpoint2026jev} & Comparator cost rose from \$0.54 to \$4.39 & The figures are \$0.56 to \$4.39 (\$0.54 is \jev{}'s own cost). The page title has changed. \\
\cite{simmons2026saidno} & ``White'' 51\% on beauty, leadership, civility, morality; 3\% ``yes'' answers & 51\% is across all eight traits (6 of 8 won); 3\% is the mean ``yes'' probability. Title and date changed. \\
\cite{qin2026zhdecisionbench} & Systematic over-confidence in Chinese; 28\% order flips & The benchmark's first release (26~September) did not test \jev{} (two Laya models and Qwen3.5-2B); the 28\% is a 25-item subset. Its second release (28~September) calls the over-confidence finding partly a small-sample artefact. \\
\cite{yurin2026confident} & Confidence formula from 76{,}807 answers & It was derived from 738{,}164 Choice answers; 76{,}807 is the two-option subset. \\
\cite{typesafe_models} & 70--500\,ms from the model documentation & The figure comes from the launch post~\cite{almeida2026introducing}; the documentation states about 100--150\,ms. \\
\pol{} text & Clauses \clause{3.4.1}, \clause{5.3.x}, \clause{3.6} & Renumbered before commit \texttt{5791ae3}, in commits \texttt{7750ceb}--\texttt{526da3d}, to \clause{5.4.1}, \clause{4.3.x}, \clause{5.6} (\cref{tab:concordance}). \\
\end{longtable}
\endgroup

\bibliographystyle{unsrtnat}
\bibliography{pol2,corpus,corpus_update,grey,background}
\end{document}
